\subsection{Overview}
\label{sec:method-overview}

% We propose \ourmodelnospace, a looped MoE language model. Following DeepSeek-V3~\cite{liu2024deepseek}, each layer pairs a Multi-head Latent Attention (MLA)~\cite{deepseekv2} sublayer with a fine-grained MoE sublayer, and the first post-embedding layer uses a dense FFN. Each sublayer adopts a sandwich normalization scheme, with pre-normalization before the sublayer and an additional RMSNorm on its output prior to the residual addition. On top of this backbone, we follow the sandwich-loop layout from prior weight-shared architectures~\cite{mcleish2025teaching,zeitoun2026hyperloop, geiping2026scaling}, organizing layers into non-shared prefix and suffix blocks around a shared loop body executed $K$ times.

% To make weight sharing across iterations expressive enough for language modeling, we introduce \textbf{Iter}ation-conditioned Token-aware \textbf{Ada}ptive \textbf{L}ayer \textbf{N}ormalization (\textbf{IterAdaLN}), which replaces the pre-normalization at every sublayer site inside the loop body and supplies per-iteration, per-token modulation. We further introduce a capacity-balancing strategy that corrects the attention–FFN active-parameter imbalance inherent to looped MoEs, and, as discussed in Section~\ref{sec:intro}, exploits MoE sparsity to align total parameters, per-token FLOPs, and sublayer active ratios with a Vanilla MoE reference (a non-looped MoE matched in capacity, detailed in Section~\ref{sec:method-budget}). This enables a clean attribution of performance gains to the loop architecture itself. Figure~\ref{fig:model_arch} illustrates the overall \ourmodel layout on the left, together with the per-iteration update inside one loop-block layer on the right.

We propose \ourmodelnospace, a looped MoE language model. Following DeepSeek-V3~\cite{liu2024deepseek}, each layer pairs a Multi-head Latent Attention (MLA)~\cite{deepseekv2} sublayer with a fine-grained MoE sublayer, and the first post-embedding layer uses a dense FFN. Each sublayer adopts a sandwich normalization scheme, with pre-normalization before the sublayer and an additional RMSNorm on its output prior to the residual addition. On top of this backbone, we follow the sandwich-loop layout from prior weight-shared architectures~\cite{mcleish2025teaching,zeitoun2026hyperloop, geiping2026scaling}, organizing layers into non-shared prefix and suffix blocks around a shared loop body executed $K$ times. Inside the loop body, we introduce IterAdaLN (Section~\ref{sec:method-tadaln}) to supply per-iteration, per-token modulation, and a capacity-balancing strategy (Section~\ref{sec:method-budget}) to align total parameters, per-token FLOPs, and active sublayer ratios with a non-looped MoE reference of matched capacity (the Vanilla MoE). Figure~\ref{fig:model_arch} illustrates the overall layout and the per-iteration update inside one loop-block layer.
\begin{figure}[t]
    \centering
    \includegraphics[width=0.95\columnwidth]{Images/model.jpg}
    \caption{Architecture of \ourmodelnospace. \textbf{Left:} the sandwich-loop layout with prefix layers, a loop block of shared layers executed $K$ times, and suffix layers. \textbf{Right:} the per-iteration update inside the Loop MoE Layer, featuring IterAdaLN at both pre-normalization sites and an asymmetric residual gate $\alpha$ on the attention branch only.
    }
    \label{fig:model_arch}
\end{figure}


%-----------------------------------------------------------------------
\subsection{Sandwich-Loop Backbone}
\label{sec:method-backbone}
%-----------------------------------------------------------------------

\ourmodel organizes its layers into three regions, namely a non-shared prefix layers $L_p$, a loop block of $L_b$ shared layers executed $K$ times, and non-shared suffix layers $L_s$. This configuration yields an effective depth of $D = L_p + K \cdot L_b + L_s$ while using only $U = L_p + L_b + L_s$ unique layers. The non-shared boundary layers provide dedicated capacity for input-embedding adaptation and output projection, while grouping the shared layers into one multi-layer block enables implicit intra-block information exchange across iterations. In particular, a refined representation produced by a later layer in iteration $k$ becomes available to an earlier layer in iteration $k\!+\!1$, allowing the block as a whole to revisit and rewrite its intermediate states.

% Within this architectural framework, the central design question is how information flows between successive iterations of the loop body. Prior sandwich-loop models~\cite{mcleish2025teaching,geiping2026scaling} commonly re-inject the prefix output additively at every iteration, that is, $h_k = f_{\mathrm{body}}(h_{k-1}) + h_{\text{pre}}$, where $h_{\text{pre}}$ denotes the output of the prefix block. This formulation has two drawbacks. First, the first iteration is ill-defined since no genuine $h_0$ exists, and is typically resolved by substituting a random vector. This injects an irrelevant signal into the loop from the very first step. Second, even after the loop has progressed, the additive pathway re-broadcasts a constant $h_{\text{pre}}$ at every step. This constant tends to dominate the residual stream over the evolving body output, suppressing the iteration-to-iteration differentiation that the loop is intended to express. We therefore remove the re-injection pathway entirely. The loop body takes $h_{\text{pre}}$ as its starting point and evolves through $K$ steps of pure recursion.

A central design question is how information flows across iterations. Prior sandwich-loop models~\cite{mcleish2025teaching,geiping2026scaling} re-inject the prefix output $h_{\text{pre}}$ additively at every step, requiring an ill-defined $h_0$ at the first iteration and allowing a constant $h_{\text{pre}}$ to dominate the residual stream, thereby suppressing iteration-to-iteration differentiation. We therefore remove the re-injection pathway so that the loop body takes $h_{\text{pre}}$ as its starting point and evolves through $K$ steps of pure recursion.
\begin{align}
\small
  h_0 &= h_{\text{pre}} \label{eq:loop-init} \\
  h_k &= f_{\mathrm{body}}(h_{k-1};\,\mathbf{c}_{k,t}), \quad k = 1, 2, \dots, K
  \label{eq:pure-recursion}
\end{align}
where $\mathbf{c}_{k,t}$ is the conditioning vector for token $t$ at step $k$ and is subsequently consumed directly by IterAdaLN (Section~\ref{sec:method-tadaln}). Each forward pass is therefore a strict recursive refinement, and IterAdaLN becomes the sole pathway for per-iteration differentiation.

%-----------------------------------------------------------------------
% \subsection{IterAdaLN: Iteration-Conditioned Token-Aware Adaptive Layer Normalization}
\subsection{IterAdaLN}
\label{sec:method-tadaln}
%-----------------------------------------------------------------------

Weight sharing constrains the model to reuse identical parameters across all steps, so the computation lacks iteration awareness without explicit conditioning. A natural remedy is to condition each iteration on a global iteration signal via Adaptive Layer Normalization (AdaLN)~\cite{peebles2023scalable}. This is sufficient for image generation, where the diffusion timestep serves as a genuinely global conditioning target. Applying the same recipe to an MoE language model backbone, however, introduces a spatial bottleneck. A purely iteration-conditioned AdaLN forces a uniform affine modulation across all sequence positions at each step $k$, whereas natural language sequences are intrinsically heterogeneous. Syntactic markers, function words, and dense semantic entities demand qualitatively different representational updates at the same iteration. With only a global signal available, all sequence-level heterogeneity must be absorbed by the weights of the shared body.

We therefore introduce \textbf{IterAdaLN}, which generates modulation parameters jointly from the iteration index and the current token state, unlocking per-iteration and per-token degrees of freedom directly within the conditioning pathway. For token position $t$ at loop iteration $k$, IterAdaLN forms branch-specific joint conditions $c^{\mathrm{attn}}_{k,t}$ and $c^{\mathrm{moe}}_{k,t}$ by fusing two streams of information. The token stream supplies local semantic context through a linear projection of the token's current state in the forward pass. For the attention branch, this state is the initial state $h_{k-1,t}$, while for the MoE branch, it is the intermediate state $m_{k,t}$ produced by the preceding attention update. The iteration stream follows the DiT-style time encoding~\cite{peebles2023scalable}, mapping the iteration index $k$ through a fixed sinusoidal positional encoding $\mathrm{PE}(k)$ and a small learnable MLP, $\mathbf{v}_k = \mathrm{MLP}_\mathrm{iter}(\mathrm{PE}(k))$. The two streams are combined by broadcast addition for each branch,

\begingroup
\small
\begin{align}
  c^{\mathrm{attn}}_{k,t} &= \mathrm{Linear}_{\mathrm{attn}}(h_{k-1,t})
              + \mathrm{MLP}_{\mathrm{iter}}\!\big(\mathrm{PE}(k)\big)
    \label{eq:c_kt_attn} \\
  c^{\mathrm{moe}}_{k,t} &= \mathrm{Linear}_{\mathrm{moe}}(m_{k,t})
              + \mathrm{MLP}_{\mathrm{iter}}\!\big(\mathrm{PE}(k)\big)
    \label{eq:c_kt_moe}
\end{align}
\endgroup
\noindent so that both conditions inherit sequence-length variation from the token term and iteration variation from the time term.

Given the corresponding condition $c$, IterAdaLN replaces standard RMSNorm~\cite{zhang2019root} at every pre-normalization site with

\begingroup
\small
\begin{equation}
  \mathrm{IterAdaLN}(h;\,c) =
     \frac{\big(1+\gamma(c)\big) \odot h}{\sqrt{\mathbb{E}[h^2]+\epsilon}}
    + \beta(c)
  \label{eq:iteradaln}
\end{equation}
\endgroup
\noindent Note that IterAdaLN employs an affine-free RMSNorm. We omit the default static learnable scale of RMSNorm, since IterAdaLN dynamically supplies a token- and iteration-conditional scale.

To provide fine-grained, independent modulation, the affine parameters are generated by two separate, zero-initialized MLPs,

\begingroup
\small
\begin{align}
  \big[\gamma^{\mathrm{attn}}_{k,t},\,\beta^{\mathrm{attn}}_{k,t},\,\alpha_{k,t}\big]
  &= \mathrm{MLP}^{\mathrm{attn}}_{\mathbf{0}}\!\big(c^{\mathrm{attn}}_{k,t}\big)
    \label{eq:mlp_attn} \\
  \big[\gamma^{\mathrm{moe}}_{k,t},\,\beta^{\mathrm{moe}}_{k,t}\big]
  &= \mathrm{MLP}^{\mathrm{moe}}_{\mathbf{0}}\!\big(c^{\mathrm{moe}}_{k,t}\big)
    \label{eq:mlp_moe}
\end{align}
\endgroup
\noindent where the subscript $\mathbf{0}$ denotes zero initialization of the output projection.

We apply the residual gate $\alpha_{k,t}$ only to the attention branch. A symmetric gate on the MoE branch would be absorbed into the routing weights and interfere with load-balancing calibration (see Appendix~\ref{app:alpha-asymmetric} for justification).


% A critical design choice concerns where the residual gate $\alpha_{k,t}$ should act. While DiT applies a symmetric gate to both sublayer outputs, we apply $\alpha_{k,t}$ exclusively to the attention branch. This design follows directly from the structure of the MoE operator. For an IterAdaLN-modulated input $\tilde h_{k,t}$, the MoE output is a routing-weighted sum over selected experts, so that a branch-level residual gate is absorbed straight into those routing weights,

% \begin{equation}
% \small
%   \alpha_{k,t}\!\cdot\!\mathrm{MoE}(\tilde h_{k,t})
%   = \!\!\sum_{i\in\mathrm{Topk=}}\!\!\big(\alpha_{k,t}\, w_{k,t,i}\big)\,
%          E_i(\tilde h_{k,t}).
%   \label{eq:gate-absorb}
% \end{equation}
% The router therefore already provides the same per-token reweighting that an MoE-branch gate would offer, so the gate adds no representational capacity. It also actively interferes with routing. The routing weights $w_{k,t,i}$ are calibrated by the load-balancing loss and the expert-bias update~\cite{liu2024deepseek} to remain in a well-behaved regime. Rescaling them by a data-dependent $\alpha_{k,t}$ perturbs that calibration along with the running statistics driving the expert-bias controller. The MLA attention sublayer, by contrast, has no analogous internal token-conditional reweighting mechanism, which makes a residual gate there strictly complementary.

We accordingly express the per-token, per-iteration layer update as the sequential composition of an attention residual sub-block and an MoE residual sub-block. The intermediate state $m_{k,t}$ is first produced by the attention block,
% \begin{equation}
% \small
% \begin{split}
%   m_{k,t} &= h_{k-1,t} + \alpha_{k,t} \, \mathcal{A} \Big( \\
%           &\quad \mathrm{IterAdaLN}(h_{k-1,t}; \, \gamma^{\mathrm{attn}}_{k,t}, \, \beta^{\mathrm{attn}}_{k,t}) \Big)
% \end{split}
% \label{eq:layer-update-attn}
% \end{equation}
\begin{equation}
\resizebox{0.95\linewidth}{!}{$
  m_{k,t} = h_{k-1,t} + \alpha_{k,t} \, \mathcal{A}\big(\mathrm{IterAdaLN}(h_{k-1,t}; \, \gamma^{\text{attn}}_{k,t}, \, \beta^{\text{attn}}_{k,t})\big)
$}
\label{eq:layer-update-attn}
\end{equation}
where $\mathcal{A}$ is the MLA attention sublayer. This updated state $m_{k,t}$ is then used both to compute the condition $c^{\mathrm{moe}}_{k,t}$ and as the input to the MoE block, yielding the final layer output,
\begin{equation}
\resizebox{0.95\linewidth}{!}{$
  h_{k,t} = m_{k,t} + \mathrm{MoE}\!\big(
         \mathrm{IterAdaLN}(m_{k,t};\, \gamma^{\mathrm{moe}}_{k,t},\, \beta^{\mathrm{moe}}_{k,t})
       \big)
$}
\label{eq:layer-update-attn}
\end{equation}
These update rules reflect the asymmetric placement of $\alpha_{k,t}$. Following AdaLN-Zero~\cite{peebles2023scalable}, the output projection of $\mathrm{MLP}_{\mathbf{0}}$ is zero-initialized, so that $\gamma_{k,t}=\beta_{k,t}=\alpha_{k,t}=0$ at initialization. The $K$-step loop therefore begins as a sequence of near-identity transformations, from which per-iteration and per-token differentiation gradually emerges during optimization. Appendix~\ref{app:loop-alpha-schedule} empirically verifies that the trained model exploits this per-iteration, per-token freedom, exhibiting a back-loaded, attention-concentrated update schedule.





%-----------------------------------------------------------------------
\subsection{Capacity Balancing Strategy}
\label{sec:method-budget}
%-----------------------------------------------------------------------

Looped MoEs exhibit a structural distortion in active capacity allocation. The dense MLA sublayer reuses the same parameters at every iteration, while a token may route to different experts across the $K$ iterations of a shared MoE layer, so that its unique active FFN parameters accumulate with $K$. The resulting ratio $\rho_{\text{loop}} = A_{\text{attn}}/A_{\text{ffn}}$ drops well below the operating point $\rho^\star$ of well-tuned non-looped MoEs, leaving attention under-parameterized relative to the FFN pool (Section~\ref{sec:ablation}). Let $N$ and $F$ denote a model's total parameter count and per-token FLOPs, with $N^\star, F^\star$ the corresponding baseline values. As noted in Section~\ref{sec:intro}, MoE sparsity decouples $N$ from $F$, so once $\rho$ is corrected we can independently restore $N$ and $F$ to $N^\star$ and $F^\star$. We achieve this alignment through a two-step procedure.

\paragraph{Measurement}
We first experimentally investigate the expected number of unique experts a token activates across the $K$ iterations of one shared layer under top-$k$ routing. The resulting $A_{\text{ffn}}$ scales sublinearly but non-trivially with $K$, while $A_{\text{attn}}$ remains constant (see Appendix~\ref{app:active-scaling}). We adopt a Vanilla MoE, a non-looped MoE, as our matched reference. We characterize it by its attention-to-FFN active ratio $\rho^\star = A^\star_{\text{attn}}/A^\star_{\text{ffn}}$, together with $N^\star$ and $F^\star$ introduced above, which jointly specify the target operating point.

\paragraph{Rebalancing and capacity restoration}
To elevate $\rho$ toward $\rho^\star$, we expand the MLA low-rank projections to increase $A_{\mathrm{attn}}$, and shrink each expert's hidden dimension to decrease $A_{\mathrm{ffn}}$, while keeping top-$k$ routing fixed. This preserves $F$ but reduces $N$. We then restore $N$ to $N^\star$ by scaling up the routed expert pool, exploiting the MoE decoupling of $N$ from $F$.